\begin{lemma}[Domain restriction and data congruence for $p$-harmonicity]
Let $1<p<2$, let $\Omega\subseteq\mathrm{Plane}$, let $u\colon\mathrm{Plane}\to\mathbb{R}$ and
$g\colon\mathrm{Plane}\to\mathrm{Plane}$, and assume that $u$ is $p$-harmonic in $\Omega$ with
weak gradient $g$.  Then:
\begin{enumerate}
\item[(a)] \textbf{Restriction.}  For every open $\Omega'\subseteq\Omega$, the function $u$ is
$p$-harmonic in $\Omega'$ with the same weak gradient $g$.
\item[(b)] \textbf{Congruence.}  For all $u'\colon\mathrm{Plane}\to\mathbb{R}$ and
$g'\colon\mathrm{Plane}\to\mathrm{Plane}$ with $u'=u$ on $\Omega$ and $g'=g$ on $\Omega$, the
function $u'$ is $p$-harmonic in $\Omega$ with weak gradient $g'$.
\end{enumerate}
\end{lemma}
\begin{proof}
We refer to the six conditions in the definition of a $p$-harmonic function with weak gradient,
\noderef{IsPHarm}, by their numbers (1)--(6), and we write $A_p$ for the $p$-flux
\noderef{PFlux}.  By condition (1) for the datum $h$, the set $\Omega$ is open, hence Lebesgue
measurable.

\medskip\noindent\emph{Step 0: a localization estimate for test functions.}
We first record a fact used in part (a).  Let $W\subseteq\mathrm{Plane}$ be open and suppose
that $u$ is continuous on $W$, that $g$ is almost everywhere strongly measurable with respect to
Lebesgue measure restricted to $W$, and that $x\mapsto\lvert g(x)\rvert^{p}$ is integrable on
every compact subset of $W$.  Let $\varphi\colon\mathrm{Plane}\to\mathbb{R}$ be of class
$C^\infty$ with compact support and $\operatorname{supp}\varphi\subseteq W$, and put
$K_0:=\operatorname{supp}\varphi$, a compact subset of $W$, and
\[
  B:=\sup_{y}\lvert D\varphi(y)\rvert<\infty,\qquad B':=\sup_{y}\lvert\varphi(y)\rvert<\infty,
\]
finite because $\varphi$ and $D\varphi$ are continuous and vanish off $K_0$ (on the open set
$\mathrm{Plane}\setminus K_0$ the function $\varphi$ is identically $0$, so all of its
derivatives vanish there).  Then each of the three functions
\[
  u\,D\varphi,\qquad \varphi\,g,\qquad \bigl\langle A_p(g),D\varphi\bigr\rangle
\]
vanishes off $K_0$ and is integrable on $W$, and its integral over $W$ equals its integral over
$K_0$.  Indeed:
\begin{itemize}
\item $W$ is open, hence measurable, by clause (2) of the background proposition on
integrability of continuous and compactly supported functions.  Since $u$ is continuous on $W$
and $D\varphi$ is continuous on $\mathrm{Plane}$, the function $u\,D\varphi$ is continuous on
$W$, hence bounded and integrable on the compact subset $K_0\subseteq W$ by clause (3) of that
proposition; it vanishes off $K_0$, so it is integrable on $W$ with
$\int_Wu\,D\varphi=\int_{K_0}u\,D\varphi$ by clause (4).
\item Taking $a=1$ in clause (4) of the background proposition on elementary properties of real
powers (legitimate since $1\le p$), $\lvert\eta\rvert\le1+\lvert\eta\rvert^{p}$ for every
vector $\eta$, so $\lvert\varphi\,g\rvert\le B'(1+\lvert g\rvert^{p})$ pointwise.  The
right-hand side is integrable on $K_0$ because $\lvert g\rvert^{p}$ is and $\lvert K_0\rvert$ is
finite.  As $\varphi\,g$ is almost everywhere strongly measurable on $W$, clause (3) of the
integrability proposition makes it integrable on $K_0$; it vanishes off $K_0$, so clause (4)
gives integrability on $W$ and $\int_W\varphi\,g=\int_{K_0}\varphi\,g$.
\item By \noderef{PFluxContinuity}, $A_p$ is continuous on $\mathrm{Plane}$ and
$\lvert A_p(\eta)\rvert=\lvert\eta\rvert^{\,p-1}$ for every $\eta$.  Hence $A_p\circ g$ is
almost everywhere strongly measurable on $W$ (composition with a continuous map preserves that,
by the background proposition on pointwise limits), and so is its inner product with the
continuous $D\varphi$.  By Cauchy--Schwarz and clause (4) of the elementary-powers proposition
with $a=p-1\in[0,p]$,
\[
  \bigl\lvert\bigl\langle A_p(g(y)),D\varphi(y)\bigr\rangle\bigr\rvert
  \le B\,\lvert g(y)\rvert^{\,p-1}\le B\bigl(1+\lvert g(y)\rvert^{p}\bigr),
\]
which is integrable on $K_0$ as before; the same two clauses give integrability on $W$ and
equality of the $W$- and $K_0$-integrals.
\end{itemize}

\medskip\noindent\emph{Part (a).}
Let $\Omega'\subseteq\Omega$ be open.  We verify the six conditions for $u$, $g$ on $\Omega'$.

(1) $\Omega'$ is open by hypothesis.

(2) $u$ is continuous on $\Omega$ by condition (2) for $h$, hence on the subset $\Omega'$, by
clause (1) of the background proposition on restriction and congruence.

(3) $g$ is almost everywhere strongly measurable for Lebesgue measure restricted to $\Omega$ by
condition (3) for $h$.  Since $\Omega'\subseteq\Omega$, clause (2) of the restriction and
congruence proposition gives the same for Lebesgue measure restricted to $\Omega'$.

(4) Let $K\subseteq\Omega'$ be compact.  Then $K\subseteq\Omega$, so condition (4) for $h$ gives
that $x\mapsto\lvert g(x)\rvert^{p}$ is integrable on $K$.

(5) Let $\varphi$ be of class $C^\infty$ with compact support and
$\operatorname{supp}\varphi\subseteq\Omega'$, and put $K_0:=\operatorname{supp}\varphi$.  Then
$K_0\subseteq\Omega'\subseteq\Omega$, so $\varphi$ is also an admissible test function for
$\Omega$, and condition (5) for $h$ gives
\[
  \int_\Omega u\,D\varphi\,dx=-\int_\Omega\varphi\,g\,dx .
\]
Step 0 applies with $W=\Omega$ (its three hypotheses are conditions (2), (3), (4) for $h$) and
also with $W=\Omega'$ (its three hypotheses are the conditions (2), (3), (4) just proved).  So
all four integrals reduce to integrals over $K_0$:
\[
  \int_{\Omega'}u\,D\varphi=\int_{K_0}u\,D\varphi=\int_\Omega u\,D\varphi,
  \qquad
  \int_{\Omega'}\varphi\,g=\int_{K_0}\varphi\,g=\int_\Omega\varphi\,g .
\]
Combining the last two displays, $\int_{\Omega'}u\,D\varphi=-\int_{\Omega'}\varphi\,g$, which
is condition (5) for $u$, $g$ on $\Omega'$.

(6) With the same $\varphi$ and $K_0$, condition (6) for $h$ gives
$\int_\Omega\langle A_p(g),D\varphi\rangle\,dx=0$, and Step 0 with $W=\Omega$ and with
$W=\Omega'$ gives
\[
  \int_{\Omega'}\bigl\langle A_p(g),D\varphi\bigr\rangle
  =\int_{K_0}\bigl\langle A_p(g),D\varphi\bigr\rangle
  =\int_\Omega\bigl\langle A_p(g),D\varphi\bigr\rangle=0 ,
\]
which is condition (6) for $u$, $g$ on $\Omega'$.  This proves (a).

\medskip\noindent\emph{Part (b).}
Let $u'=u$ on $\Omega$ and $g'=g$ on $\Omega$.  We verify the six conditions for $u'$, $g'$ on
$\Omega$.  The point throughout is that every condition is either an assertion about points of
$\Omega$ or an integral over $\Omega$, so it only sees the restrictions of the data to $\Omega$.

(1) $\Omega$ is open, by condition (1) for $h$; this condition does not involve $u$ or $g$.

(2) $u$ is continuous on $\Omega$ by condition (2) for $h$, and $u'=u$ on $\Omega$, so $u'$ is
continuous on $\Omega$ by clause (1) of the restriction and congruence proposition.

(3) $g$ is almost everywhere strongly measurable for Lebesgue measure restricted to $\Omega$ by
condition (3) for $h$.  Since $g'=g$ on the measurable set $\Omega$, clause (3) of the
restriction and congruence proposition gives first that $g'=g$ almost everywhere for that
measure, and then that $g'$ is almost everywhere strongly measurable for it.

(4) Let $K\subseteq\Omega$ be compact.  For $x\in K$ we have $g'(x)=g(x)$ and hence
$\lvert g'(x)\rvert^{p}=\lvert g(x)\rvert^{p}$, so the two integrands agree on $K$.  By
condition (4) for $h$ the second is integrable on $K$, so by clause (5) of the restriction and
congruence proposition the first is as well.

(5) Let $\varphi$ be of class $C^\infty$ with compact support and
$\operatorname{supp}\varphi\subseteq\Omega$.  The integrands $u'\,D\varphi$ and $u\,D\varphi$
agree at every point of $\Omega$, and so do $\varphi\,g'$ and $\varphi\,g$.  By clause (5) of
the restriction and congruence proposition their integrals over $\Omega$ agree:
\[
  \int_\Omega u'\,D\varphi=\int_\Omega u\,D\varphi,\qquad
  \int_\Omega\varphi\,g'=\int_\Omega\varphi\,g .
\]
Condition (5) for $h$ equates the right-hand sides up to sign, so
$\int_\Omega u'\,D\varphi=-\int_\Omega\varphi\,g'$, which is condition (5) for $u'$, $g'$.

(6) With the same $\varphi$: for $x\in\Omega$ we have $g'(x)=g(x)$ and therefore
$A_p(g'(x))=A_p(g(x))$, so the integrands $\langle A_p(g'),D\varphi\rangle$ and
$\langle A_p(g),D\varphi\rangle$ agree at every point of $\Omega$.  By clause (5) of the
restriction and congruence proposition their integrals over $\Omega$ agree, and condition (6)
for $h$ makes the second $0$.  Hence
$\int_\Omega\langle A_p(g'),D\varphi\rangle\,dx=0$, which is condition (6) for $u'$, $g'$.

\medskip\noindent Both parts are proved.
\end{proof}
